\documentclass[12pt,a4paper]{article}

\usepackage[utf8]{inputenc}
\usepackage[T1]{fontenc}
\usepackage[brazil]{babel}
\usepackage[a4paper,margin=2.5cm]{geometry}
\usepackage{listings}
\usepackage{xcolor}
\usepackage{hyperref}
\usepackage{enumitem}
\usepackage{titlesec}
\usepackage{parskip}
\usepackage{graphicx}
\usepackage{array}
\usepackage{tabularx}
\usepackage{seqsplit}
\usepackage{colortbl}

% Helper para texto monoespacado com quebra automatica (evita overflow em tabelas).
% Importante: escape caracteres especiais no input -- _ como \_ , < como \textless ,
% > como \textgreater . O seqsplit quebra a string caractere a caractere.
\newcommand{\ttwrap}[1]{{\ttfamily\seqsplit{#1}}}

\hypersetup{
    colorlinks=true,
    linkcolor=blue!60!black,
    urlcolor=blue!60!black,
    citecolor=blue!60!black,
}

% --- Configuração para listings (TypeScript / Python) ---
\definecolor{codegreen}{rgb}{0.0,0.5,0.0}
\definecolor{codegray}{rgb}{0.4,0.4,0.4}
\definecolor{codeblue}{rgb}{0.13,0.29,0.53}
\definecolor{codepurple}{rgb}{0.58,0.0,0.83}
\definecolor{backcolour}{rgb}{0.96,0.96,0.96}

\lstdefinestyle{codestyle}{
    backgroundcolor=\color{backcolour},
    commentstyle=\color{codegreen},
    keywordstyle=\color{codeblue}\bfseries,
    numberstyle=\tiny\color{codegray},
    stringstyle=\color{codepurple},
    basicstyle=\ttfamily\footnotesize,
    breakatwhitespace=false,
    breaklines=true,
    captionpos=b,
    keepspaces=true,
    numbers=left,
    numbersep=6pt,
    showspaces=false,
    showstringspaces=false,
    showtabs=false,
    tabsize=2,
    frame=single,
    rulecolor=\color{gray!40},
    framesep=4pt,
}
\lstset{style=codestyle}

% Linguagem TypeScript "manual" (não nativa no listings)
\lstdefinelanguage{TypeScript}{
  keywords={break, case, catch, class, const, continue, debugger, default,
            delete, do, else, enum, export, extends, false, finally, for,
            function, if, implements, import, in, instanceof, interface,
            let, new, null, package, private, protected, public, return,
            static, super, switch, this, throw, true, try, typeof, var,
            void, while, with, yield, async, await, of, as, from, type,
            describe, it, expect, jest, beforeEach, afterEach},
  keywordstyle=\color{codeblue}\bfseries,
  ndkeywords={Number, String, Boolean, Promise, Array, Request, Response, void},
  ndkeywordstyle=\color{codeblue!70!black}\bfseries,
  identifierstyle=\color{black},
  sensitive=true,
  comment=[l]{//},
  morecomment=[s]{/*}{*/},
  commentstyle=\color{codegreen}\ttfamily,
  stringstyle=\color{codepurple}\ttfamily,
  morestring=[b]',
  morestring=[b]"
}

\titleformat{\section}{\Large\bfseries\color{codeblue}}{\thesection}{1em}{}
\titleformat{\subsection}{\large\bfseries\color{codeblue!80!black}}{\thesubsection}{1em}{}

% ---------- Cabeçalho institucional "UnB | FGA" estilizado ----------
\definecolor{unbgreen}{rgb}{0.0,0.39,0.20}   % verde institucional
\definecolor{unbblue}{rgb}{0.0,0.20,0.40}    % azul institucional

\newcommand{\unbfgaheader}{%
  \begin{center}
    \setlength{\fboxsep}{6pt}
    \setlength{\fboxrule}{1.2pt}
    {\Large
      \fcolorbox{unbgreen}{white}{\color{unbgreen}\bfseries UnB}%
      \hspace{0.3em}%
      \fcolorbox{unbblue}{unbblue}{\color{white}\bfseries\hspace{0.2em}FGA\hspace{0.2em}}%
    }\\[0.3em]
    {\small\itshape Universidade de Brasília\,\textbar\, Faculdade do Gama}\\
    {\small\itshape Engenharia de Software}\\
    {\small\itshape Disciplina FGA0314 -- Testes de Software}
  \end{center}
  \vspace{0.4em}
  \noindent\textcolor{unbgreen}{\rule{\linewidth}{1.5pt}}
  \vspace{0.6em}
}

\begin{document}

% Cabeçalho institucional substitui o \maketitle padrão
\unbfgaheader

\begin{center}
    {\LARGE\bfseries\color{unbblue}
      Relatório de Auditoria de Testes -- Módulo 3}\\[0.4em]
    {\large
      Atividades 1, 2 e 3 -- Auditoria dos Testes de Unidade do Projeto NoFluxo UnB}\\[1.2em]
\end{center}

% ---------- Bloco de integrantes ----------
\begin{center}
\footnotesize
\begin{tabularx}{0.85\textwidth}{|>{\raggedright\arraybackslash}X|>{\raggedright\arraybackslash}p{3.5cm}|}
\hline
\rowcolor{backcolour}\textbf{Integrante} & \textbf{Matrícula} \\ \hline
Enzo Lopes Ferreira                    & 232002000 \\ \hline
Vinícius de Jesus Bessa Fernandes      & 222006490 \\ \hline
Vitor Marconi Trancoso Albuquerque     & 222006202 \\ \hline
\end{tabularx}
\end{center}

\begin{center}
\small Brasília, maio de 2026
\end{center}

\vspace{1em}

\section*{Contexto da auditoria}

Esta equipe atua como \textbf{colaboradora externa} auditando os testes
do projeto \textbf{NoFluxo UnB}. A equipe \emph{não} desenvolveu o
código de produção: o trabalho aqui descrito é de \emph{análise} dos
testes existentes (Atividades 1 e 2) e de \emph{proposta} de uma
reescrita exemplar (Atividade 3).

O projeto auditado é composto por:

\begin{itemize}[leftmargin=1.5em]
    \item Um \textbf{backend Node.js/TypeScript} (\texttt{no\_fluxo\_backend}) que expõe controllers Express e usa Supabase como banco;
    \item Um \textbf{serviço Python} de IA (\texttt{ai\_agent}) baseado em Flask, integrado ao RAGFlow;
    \item Scripts \textbf{Python} de scraping/coleta de dados do SIGAA;
    \item Frontends em \textbf{Flutter} e \textbf{Svelte}.
\end{itemize}

A auditoria concentrou-se em duas áreas com testes \emph{de unidade}:
\texttt{no\_fluxo\_backend/tests-ts/} (Jest + TypeScript) e
\texttt{tests-python/} (Pytest + unittest). Frontends e scripts
sem testes de unidade foram excluídos do escopo.

% =====================================================================
\section{Atividade 1 -- Anatomia dos Testes do Projeto}
% =====================================================================

O objetivo desta atividade é mapear, do ponto de vista da auditoria,
como o projeto NoFluxo organiza seus testes de unidade. A equipe
analisou as quatro dimensões pedidas no enunciado: convenção de nomes,
estrutura AAA, uso de mocks/stubs e dependências externas.

\subsection{Estrutura geral dos testes}

\begin{center}
\footnotesize
\begin{tabularx}{\textwidth}{|>{\raggedright\arraybackslash}p{3.4cm}|>{\raggedright\arraybackslash}p{3.6cm}|>{\raggedright\arraybackslash}X|}
\hline
\textbf{Área} & \textbf{Framework} & \textbf{Localização} \\ \hline
Backend TS (Express) & Jest 29.7.0 + ts-jest & \ttwrap{no\_fluxo\_backend/tests-ts/*.test.ts} \\ \hline
Serviço Python (Flask) & Pytest + pytest-mock & \ttwrap{tests-python/test\_*.py} \\ \hline
Scrapers Python & unittest + mock & \ttwrap{tests-python/test\_extrair\_*.py} \\ \hline
\end{tabularx}
\end{center}

\subsection{Convenção de nomes}

\textbf{No backend TypeScript} o projeto segue a convenção do Jest com
\texttt{describe} agrupando por controller e \texttt{it} descrevendo o
comportamento em inglês na forma \emph{``should + verbo + condição''}:

\begin{lstlisting}[language=TypeScript,caption={\texttt{cursos\_controller.test.ts} -- linhas 22, 42, 74, 90}]
// 'describe' agrupa os testes de um mesmo controller (boa pratica)
describe('CursosController', () => {
  // Cada 'it' descreve UM comportamento -- formato "should + acao + condicao"
  it('should return 200 with cursos data when successful', async () => { ... });
  it('should return 500 when cursos query fails', async () => { ... });
  it('should return 500 when creditos_por_curso query fails', async () => { ... });
});
// >>> Observacao da auditoria: nomes em ingles, foco em comportamento -- OK
\end{lstlisting}

\textbf{Nos testes Python} a convenção é o prefixo \texttt{test\_}
seguido por uma descrição em \textbf{português} usando \texttt{snake\_case}:

\begin{lstlisting}[language=Python,caption={\texttt{tests-python/test\_app.py} -- linha 55}]
# Convencao Python: prefixo test_ + descricao em portugues, snake_case
# 'client' e 'mocker' sao fixtures injetadas pelo pytest
def test_analisar_materia_endpoint_sucesso(client, mocker):
    """Testa o endpoint /assistente em um cenário de sucesso..."""
# >>> Observacao da auditoria: o nome cita o endpoint e o cenario,
#     mas nao explicita o RESULTADO esperado (slide 24 da aula)
\end{lstlisting}

\textbf{Avaliação da auditoria:} a convenção do backend TS está
alinhada às boas práticas (foco em comportamento). Já a Python está
parcialmente alinhada -- os nomes descrevem comportamento, mas em
alguns testes herdados de \texttt{unittest} usam-se nomes muito
genéricos como \texttt{test\_remover\_acentos} ou
\texttt{test\_coleta\_dados}, que focam em \emph{nome do método} e
não no comportamento esperado, exatamente o problema descrito no
slide 23 da aula.

\subsection{Estrutura AAA (Arrange / Act / Assert)}

A equipe identificou que os testes do backend TS seguem AAA de forma
\textbf{implícita}: o padrão está presente, mas \emph{sem} comentários
demarcando os blocos. Exemplo extraído do projeto:

\begin{lstlisting}[language=TypeScript,caption={\texttt{cursos\_controller.test.ts} -- linhas 42--72 (AAA implícito)}]
it('should return 200 with cursos data when successful', async () => {
    // ===== ARRANGE =====
    // (no projeto auditado NAO ha esses comentarios -- estamos
    //  destacando os blocos AAA para fins didaticos)
    const mockCursosData = [ { id_curso: 1, nome_curso: 'Administracao' }, ... ];

    // Configura DUAS respostas sequenciais do banco mockado:
    //  1a chamada -> dados de cursos
    //  2a chamada -> dados de creditos por curso
    (SupabaseWrapper.get().from('cursos').select as jest.Mock)
      .mockResolvedValueOnce({ data: mockCursosData,    error: null })
      .mockResolvedValueOnce({ data: mockCreditosData,  error: null });

    // ===== ACT =====
    // UMA unica chamada ao handler do controller -- boa pratica
    const handler = CursosController.routes["all-cursos"].value;
    await handler(mockRequest as Request, mockResponse as Response);

    // ===== ASSERT =====
    // Verifica status HTTP + corpo da resposta
    expect(statusSpy).toHaveBeenCalledWith(200);
    expect(jsonSpy).toHaveBeenCalledWith([ ... ]);
});
\end{lstlisting}

Nos testes Python a equipe observou que o padrão AAA também aparece,
mas é frequentemente \textbf{misturado}: o setup do cliente Flask é
uma \emph{fixture}, mocks podem aparecer espalhados, e em alguns
casos há vários blocos de \texttt{Act + Assert} no mesmo teste.

\subsection{Uso de mocks/stubs}

\textbf{Backend TS:} o projeto usa exclusivamente \emph{Jest mocks}
(\texttt{jest.fn()}, \texttt{jest.mock()}, \texttt{jest.spyOn()}).
A estratégia adotada é mockar o módulo \texttt{SupabaseWrapper} no
topo de cada arquivo de teste, evitando qualquer chamada real ao
banco:

\begin{lstlisting}[language=TypeScript,caption={\texttt{cursos\_controller.test.ts} -- linhas 6--20 (mock do Supabase)}]
// Substitui o modulo 'supabase_wrapper' inteiro por um dublê
// (NENHUMA chamada real ao banco acontece nos testes)
jest.mock('../src/supabase_wrapper', () => {
  // Cada nivel da cadeia .from(...).select(...) precisa ser mockado
  const mockSelect = jest.fn();
  const mockFrom   = jest.fn(() => ({ select: mockSelect }));
  const mockGet    = jest.fn(() => ({ from:   mockFrom   }));

  // Retorna o objeto que substitui o modulo original
  return { SupabaseWrapper: { get: mockGet } };
});
// >>> Auditoria: bom isolamento, MAS o teste fica acoplado a forma
//     exata como o controller chama o Supabase (quebra em refactor)
\end{lstlisting}

\textbf{Python:} o projeto usa \texttt{unittest.mock} (\texttt{@patch},
\texttt{MagicMock}, \texttt{mock\_open}) e o auxiliar
\texttt{pytest-mock} (\texttt{mocker.patch}). Exemplo encontrado em
\texttt{test\_app.py}:

\begin{lstlisting}[language=Python,caption={\texttt{tests-python/test\_app.py} -- linhas 60--63}]
# 'mocker.patch' substitui RagflowClient (chamada HTTP externa)
# por um dublê controlado pelo teste
mock_ragflow_client = mocker.patch(
    'no_fluxo_backend.ai_agent.app.RagflowClient'
).return_value

# Define respostas estaticas para os metodos usados no fluxo
mock_ragflow_client.start_session.return_value = "sessao_mock_123"
mock_ragflow_client.analyze_materia.return_value = {
    "code": 0, "data": "dados_mock"
}
# >>> Auditoria: dublê bem aplicado -- isola API externa do RAGFlow
\end{lstlisting}

\subsection{Dependências externas}

A tabela abaixo resume as dependências externas \emph{detectadas pela
auditoria} e como elas são tratadas nos testes existentes do projeto:

\begin{center}
\footnotesize
\begin{tabularx}{\textwidth}{|>{\raggedright\arraybackslash}p{3.5cm}|>{\raggedright\arraybackslash}X|>{\raggedright\arraybackslash}X|}
\hline
\textbf{Dependência} & \textbf{Onde aparece} & \textbf{Tratamento nos testes} \\ \hline
Banco Supabase & Backend TS, todos os controllers & Mockado via \ttwrap{jest.mock} \\ \hline
\ttwrap{Request}/\ttwrap{Response} Express & Backend TS & Mock manual via \ttwrap{Partial\textless{}Request\textgreater{}} \\ \hline
RAGFlow API (HTTP) & \ttwrap{ai\_agent/app.py} & \ttwrap{mocker.patch} \\ \hline
SIGAA (HTTP scraping) & \ttwrap{coleta\_dados/scraping/} & \ttwrap{@patch('requests.Session')} \\ \hline
Sistema de arquivos & Salvar/ler JSON e CSV & \ttwrap{mock\_open}, \ttwrap{patch('builtins.open')} \\ \hline
Variáveis de ambiente & \ttwrap{config.py} & \ttwrap{monkeypatch} (pytest) \\ \hline
Relógio/data atual & Não há testes que dependam disso & N/A \\ \hline
\end{tabularx}
\end{center}

\textbf{Conclusão da auditoria (Atividade 1):} o projeto \emph{já
adota parcialmente} as boas práticas vistas em aula -- nomeação por
comportamento, AAA, e mocks para isolar dependências externas. O
backend TS é a área mais alinhada com as recomendações; a parte
Python apresenta mais inconsistências (nomes genéricos, AAA
misturado) e alguns testes com setups muito longos, o que será
explorado nas Atividades 2 e 3.

% =====================================================================
\section{Atividade 2 -- Encontrar um ``Teste Ruim''}
% =====================================================================

\subsection{Teste escolhido pela auditoria}

A equipe selecionou o teste
\texttt{``should return 200 with database test results when successful''}
do bloco \texttt{describe('banco')}, localizado no arquivo
\texttt{no\_fluxo\_backend/tests-ts/testes\_controller.test.ts}
(linhas 67--130). Trata-se de um teste já presente no projeto que
acumula vários sintomas de ``teste ruim'' descritos nos slides 22--25
da aula.

\subsection{Código original (extraído do projeto auditado)}

\begin{lstlisting}[language=TypeScript,caption={\texttt{testes\_controller.test.ts} -- linhas 67--130 (teste original ``ruim'') -- comentarios da auditoria em destaque}]
describe('banco', () => {
    // [PROBLEMA 6] Nome generico: nao diz QUAL cenario gera o 200
    it('should return 200 with database test results when successful', async () => {

      // [PROBLEMA 5] Dados de teste excessivos -- nada disso eh
      // verificado nos asserts; aumenta carga cognitiva
      const mockCursosData = [
        { nome_curso: 'Administracao' },
        { nome_curso: 'Engenharia' }
      ];

      const mockMateriasData = [
        {
          nome_curso: 'Administracao',
          materias_por_curso: [
            {
              id_materia: 1,
              nivel: 1,
              materias: {
                id_materia: 1,
                nome_materia: 'Matematica',
                codigo_materia: 'MAT001'
              }
            },
            {
              id_materia: 2,
              nivel: 2,
              materias: {
                id_materia: 2,
                nome_materia: 'Portugues',
                codigo_materia: 'PORT001'
              }
            }
          ]
        }
      ];

      // [PROBLEMA 4] Os mocks abaixo modelam a forma EXATA como o
      // controller chama o Supabase (.from -> .select -> .limit).
      // Qualquer refactor interno quebra o teste.
      const mockSelectChain1 = jest.fn().mockResolvedValueOnce({
        data: mockCursosData,
        error: null,
      });

      const mockSelectChain2 = jest.fn().mockResolvedValueOnce({
        data: mockMateriasData,
        error: null,
      });

      const mockLimitChain = jest.fn().mockReturnValue({
        select: mockSelectChain2,
      });

      // [PROBLEMA 3] Arrange gigante e sem separacao; AAA implícito
      mockFrom.mockReturnValueOnce({
        select: mockSelectChain1,
      }).mockReturnValueOnce({
        select: jest.fn().mockReturnValue({
          limit: mockLimitChain,
        }),
      });

      // ----- aqui acaba o Arrange e comeca Act + Assert (sem demarcacao) -----
      const handler = TestesController.routes["banco"].value;
      await handler(mockRequest as Request, mockResponse as Response);

      // [PROBLEMA 2] TRES asserts verificando coisas distintas:
      // status HTTP, identificador do teste, status logico.
      // Quando falha, nao se sabe QUAL comportamento foi violado.
      expect(statusSpy).toHaveBeenCalledWith(200);
      const responseData = jsonSpy.mock.calls[0][0];
      expect(responseData.teste).toBe('conexao_banco');
      expect(responseData.status).toBe('sucesso');
      // [PROBLEMA 7] Nao ha cobertura de valores limite (lista vazia,
      // erro de query, 1 unico curso) -- so o caminho feliz
    });
});
// [PROBLEMA 1] O teste inteiro tem ~65 linhas (>90% sao setup de mock)
// [PROBLEMA 8] mockMateriasData se repete em describe('curso') e
// describe('casamento') -- duplicacao indica falta de test data builder
\end{lstlisting}

\subsection{Problemas identificados pela auditoria}

\begin{enumerate}[leftmargin=1.5em]

\item \textbf{Difícil de entender (mais de 60 linhas).} O teste tem
\textasciitilde{}65 linhas. A maior parte é setup de mock; o
comportamento testado real (uma chamada e três asserts) está enterrado
ao final.

\item \textbf{Múltiplos asserts verificando coisas diferentes.}
\texttt{expect(statusSpy).toHaveBeenCalledWith(200)},
\texttt{expect(responseData.teste).toBe(...)} e
\texttt{expect(responseData.status).toBe(...)} testam
\emph{três} aspectos distintos. Quando o teste falha, não fica claro
\emph{qual} comportamento foi violado -- o nome do teste só promete
``return 200''.

\item \textbf{Mistura Arrange/Act/Assert.} Não há separação clara entre
os três blocos. O Arrange é gigante, sem comentários, e a inicialização
\texttt{mockFrom.mockReturnValueOnce(...)} é encadeada com lógica de
``primeira chamada/segunda chamada'' que esconde a regra de negócio
testada.

\item \textbf{Depende fortemente da implementação interna.} Os mocks
modelam a forma exata como o controller chama o Supabase
(\texttt{from -> select -> limit -> select}). Qualquer refactor
interno que mantenha o mesmo comportamento externo \emph{quebra} o
teste -- problema clássico discutido no slide 23 da aula
(``foca na implementação'').

\item \textbf{Dados de teste excessivos e irrelevantes.} O teste cria
duas matérias completas (\texttt{Matematica}, \texttt{Portugues}) com
níveis e códigos -- nada disso é verificado nos asserts. Isso
aumenta a carga cognitiva sem agregar valor.

\item \textbf{Nome genérico, não comunica o cenário.}
\texttt{``should return 200 with database test results when successful''}
diz \emph{o que} retorna, mas não \emph{em que cenário} (cursos
encontrados? base vazia? quantos itens?). Pelo padrão da aula
(slide 24), o nome deveria descrever o comportamento de forma
inequívoca.

\item \textbf{Não cobre valores limite.} Não há testes para:
lista de cursos vazia, erro na segunda query (\texttt{materias\_por\_curso}),
ou um único curso retornado. Apenas o caminho ``feliz'' está coberto.

\item \textbf{Setup duplicado entre testes do mesmo \texttt{describe}.}
O mesmo \texttt{mockMateriasData} aparece praticamente igual em
\texttt{describe('curso')} e \texttt{describe('casamento')}, indicando
falta de \emph{test data builders}.

\end{enumerate}

% =====================================================================
\section{Atividade 3 -- Reescrita do Teste}
% =====================================================================

A equipe propõe, a seguir, uma reescrita do teste auditado.
A proposta aplica AAA explícito, melhora a nomenclatura, separa o
teste em casos focados (\emph{um act, um comportamento}) e adiciona
casos de valores limite (lista vazia, erro de banco), conforme pedido
no enunciado. Esta seção representa a contribuição direta da equipe
auditora ao projeto -- e é a única do relatório escrita em primeira
pessoa.

\subsection{Test data builder e helper de mock}

Primeiro, extraímos os dados e o setup repetido para um helper.
Isso reduz o ruído nos testes propostos e remove duplicação que a
auditoria identificou no arquivo original.

\begin{lstlisting}[language=TypeScript,caption={Helpers propostos pela equipe -- reduzem setup repetido}]
// =====================================================================
// HELPERS DE TESTE -- declarados no topo do arquivo
// Objetivo: tirar o ruido de mock dos casos de teste e remover
// duplicacao identificada pela auditoria (PROBLEMA 8).
// =====================================================================

// [BUILDER] Cria um curso minimo, so com o que importa pro teste
const aCurso = (nome: string) => ({ nome_curso: nome });

// [HELPER] Encapsula o mock da 1a query (lista de cursos)
//   data  -> resposta de sucesso
//   error -> caso de erro de banco (default: null = sucesso)
const stubCursosQuery = (data: any, error: any = null) => {
    const select = jest.fn().mockResolvedValueOnce({ data, error });
    mockFrom.mockReturnValueOnce({ select });
};

// [HELPER] Encapsula o mock da 2a query (materias_por_curso)
// Esconde o encadeamento .from -> .select -> .limit -> .select
// dos casos de teste -- assim, mudancas internas afetam apenas
// este helper, nao todos os testes.
const stubMateriasPorCursoQuery = (data: any, error: any = null) => {
    const select = jest.fn().mockResolvedValueOnce({ data, error });
    const limit  = jest.fn().mockReturnValue({ select });
    mockFrom.mockReturnValueOnce({
        select: jest.fn().mockReturnValue({ limit }),
    });
};
\end{lstlisting}

\subsection{Testes reescritos pela equipe auditora}

Na proposta da equipe, cada teste foca em \emph{um} comportamento,
segue AAA explícito e cobre um cenário limite distinto.

\begin{lstlisting}[language=TypeScript,caption={Versão reescrita aplicando AAA, melhor nomenclatura e cobertura de valores limite}]
// =====================================================================
// SUITE REESCRITA PELA EQUIPE AUDITORA
// Os 5 testes abaixo SUBSTITUEM o teste unico original (~65 linhas)
// e cobrem caminho feliz + 4 cenarios de borda.
// =====================================================================
describe('TestesController.banco (teste de conexao com o banco)', () => {

    // -----------------------------------------------------------------
    // CASO 1 -- caminho feliz com varios cursos
    // -----------------------------------------------------------------
    it('Banco_com_cursos_e_materias_retorna_status_sucesso', async () => {
        // Arrange -- prepara 2 cursos e 1 entrada de materias
        stubCursosQuery([aCurso('Administracao'), aCurso('Engenharia')]);
        stubMateriasPorCursoQuery([{
            nome_curso: 'Administracao',
            materias_por_curso: [],
        }]);

        // Act -- UMA unica chamada ao handler (slide 13)
        await TestesController.routes["banco"].value(
            mockRequest as Request,
            mockResponse as Response,
        );

        // Assert -- objectContaining: nao quebra se a resposta ganhar
        // novos campos no futuro (asserts focados em comportamento)
        expect(statusSpy).toHaveBeenCalledWith(200);
        expect(jsonSpy).toHaveBeenCalledWith(
            expect.objectContaining({
                teste:  'conexao_banco',
                status: 'sucesso',
            }),
        );
    });

    // -----------------------------------------------------------------
    // CASO 2 -- valor limite: ZERO cursos cadastrados
    // -----------------------------------------------------------------
    it('Banco_sem_cursos_cadastrados_retorna_status_sucesso_com_lista_vazia', async () => {
        // Arrange -- limite inferior: lista vazia
        stubCursosQuery([]);
        stubMateriasPorCursoQuery([]);

        // Act
        await TestesController.routes["banco"].value(
            mockRequest as Request,
            mockResponse as Response,
        );

        // Assert -- mesmo sem cursos, conexao deve responder sucesso
        expect(statusSpy).toHaveBeenCalledWith(200);
        expect(jsonSpy).toHaveBeenCalledWith(
            expect.objectContaining({
                teste:  'conexao_banco',
                status: 'sucesso',
            }),
        );
    });

    // -----------------------------------------------------------------
    // CASO 3 -- valor limite: EXATAMENTE 1 curso
    // -----------------------------------------------------------------
    it('Banco_com_um_unico_curso_retorna_status_sucesso', async () => {
        // Arrange -- limite minimo de uma colecao nao vazia
        stubCursosQuery([aCurso('Engenharia')]);
        stubMateriasPorCursoQuery([{
            nome_curso: 'Engenharia',
            materias_por_curso: [],
        }]);

        // Act
        await TestesController.routes["banco"].value(
            mockRequest as Request,
            mockResponse as Response,
        );

        // Assert -- foco em: handler nao falha com colecao de 1 item
        expect(statusSpy).toHaveBeenCalledWith(200);
    });

    // -----------------------------------------------------------------
    // CASO 4 -- caminho de ERRO: 1a query (cursos) falha
    // -----------------------------------------------------------------
    it('Banco_quando_query_de_cursos_falha_retorna_500_com_mensagem_de_erro', async () => {
        // Arrange -- simula erro retornado pelo Supabase
        const erro = { message: 'Database error' };
        stubCursosQuery(null, erro);

        // Act
        await TestesController.routes["banco"].value(
            mockRequest as Request,
            mockResponse as Response,
        );

        // Assert -- payload de erro deve conter mensagem original
        expect(statusSpy).toHaveBeenCalledWith(500);
        expect(jsonSpy).toHaveBeenCalledWith({
            teste:  'conexao_banco',
            status: 'erro',
            erro:   erro.message,
        });
    });

    // -----------------------------------------------------------------
    // CASO 5 -- caminho de ERRO: 2a query (materias) falha
    //   (este caso NAO existia no teste original)
    // -----------------------------------------------------------------
    it('Banco_quando_query_de_materias_falha_retorna_500_com_mensagem_de_erro', async () => {
        // Arrange -- 1a query OK, 2a query erro
        stubCursosQuery([aCurso('Administracao')]);
        stubMateriasPorCursoQuery(null, { message: 'materias error' });

        // Act
        await TestesController.routes["banco"].value(
            mockRequest as Request,
            mockResponse as Response,
        );

        // Assert -- basta verificar status 500 + flag de erro logica
        expect(statusSpy).toHaveBeenCalledWith(500);
        expect(jsonSpy).toHaveBeenCalledWith(
            expect.objectContaining({ status: 'erro' }),
        );
    });
});
\end{lstlisting}

\subsection{Melhorias propostas pela equipe auditora}

\begin{enumerate}[leftmargin=1.5em]

\item \textbf{AAA explícito.} Cada teste proposto contém comentários
\texttt{// Arrange}, \texttt{// Act} e \texttt{// Assert} demarcando os
blocos -- aumenta a legibilidade (slide 18 da aula).

\item \textbf{Nomes que descrevem comportamento.} Em vez de
\texttt{should return 200 with database test results when successful},
a equipe propõe
\texttt{Banco\_com\_cursos\_e\_materias\_retorna\_status\_sucesso}.
O nome descreve \emph{o cenário} (com cursos e matérias) e
\emph{o resultado esperado} (status sucesso), seguindo a recomendação
do slide 24 (``Use frases simples e expressivas'') e do slide 25
(separar palavras com \emph{underscore}).

\item \textbf{Cobertura de valores limite.} O teste original do
projeto cobria apenas o caminho feliz. A versão proposta pela equipe
adiciona:
\begin{itemize}[leftmargin=1.5em]
    \item Lista de cursos \emph{vazia} (zero elementos);
    \item \emph{Um único} curso (limite mínimo de uma coleção não vazia);
    \item Falha na primeira query;
    \item Falha na segunda query (materias\_por\_curso).
\end{itemize}

\item \textbf{Um Act por teste.} Cada teste proposto tem exatamente
uma chamada ao handler -- prática recomendada no slide 13 (``Um bom
teste normalmente possui apenas UM Act'').

\item \textbf{Asserts focados em comportamento, não em estrutura.}
\texttt{expect.objectContaining(\{ status: 'sucesso' \})} verifica
apenas o que importa para aquele teste, em vez de exigir o JSON
completo.  Isso reduz a fragilidade quando campos auxiliares são
adicionados/removidos da resposta.

\item \textbf{Helpers reduzem ruído de setup.} A equipe introduziu
\texttt{stubCursosQuery} e \texttt{stubMateriasPorCursoQuery} para
encapsular a forma como os mocks são encadeados, permitindo reuso
entre testes e mantendo o foco do caso de teste no \emph{cenário},
não no detalhe de implementação.

\item \textbf{Independência total.} Cada teste proposto prepara seus
próprios mocks (graças ao \texttt{afterEach} com
\texttt{jest.clearAllMocks()}) e não depende da ordem de execução --
atende ao requisito do slide 19 (``um teste deve poder executar
sozinho, em qualquer ordem'').

\end{enumerate}

% =====================================================================
\section*{Conclusão da auditoria}
% =====================================================================

A auditoria realizada pela equipe colaboradora mostrou que o projeto
NoFluxo já adota \emph{parcialmente} as boas práticas vistas no
Módulo 3 -- isolamento via mocks, AAA implícito e nomes razoáveis no
backend TypeScript. Os pontos fracos detectados foram:
(i) testes longos com setup excessivo;
(ii) acoplamento dos mocks à implementação interna do Supabase
(testes que dependem do encadeamento exato \texttt{from -> select -> limit});
(iii) ausência de testes para valores limite e caminhos de erro nos
testes do \texttt{TestesController}; e
(iv) nomes que descrevem o método/retorno em vez do comportamento.

A reescrita proposta na Atividade 3 ilustra como aplicar diretamente
as recomendações da aula (AAA explícito, nomes expressivos, um Act
por teste, cobertura de limites) \emph{sem} alterar o código de
produção do NoFluxo -- ganhando legibilidade e cobertura ao mesmo
tempo. A equipe recomenda que essa abordagem seja estendida para os
demais arquivos de teste do projeto em ciclos futuros de auditoria.

\end{document}
